\par\medskip\noindent\textbf{Cómo leer la salida.}
{\small\setlength{\tabcolsep}{4pt}\renewcommand{\arraystretch}{1.15}
\begin{xltabular}{\linewidth}{@{}>{\raggedright\arraybackslash}p{0.24\linewidth}>{\raggedleft\arraybackslash}p{0.15\linewidth}>{\raggedright\arraybackslash}X@{}}
\toprule Indicador & Valor & Interpretación \\ \midrule \endfirsthead
\toprule Indicador & Valor & Interpretación \\ \midrule \endhead
\bottomrule \endlastfoot
Días de crisis / resto & 504 / 4526 & Crisis financiera de jul-2007 a jun-2009 frente al resto. \\
Pearson, $\tau$ & 0,953 y 0,799 / 0,917 y 0,740 & La dependencia aumenta en la crisis con ambas medidas. \\
$\hat\lambda$ al 5\,\% & 0,833 y 0,794 / 0,685 y 0,681 & La dependencia en las colas sube mucho en la crisis, sobre todo en la cola inferior. \\
$t$: $\nu$ crisis / resto & 4,33 / 5,37 & Colas de la cópula más pesadas en la crisis. \\
$\lambda$ de la $t$ & 0,727 / 0,619 & La probabilidad de caídas extremas conjuntas es mayor en la crisis. \\
$\tau$ móvil: mín. & 0,585 (2000-08) & Menor dependencia en la burbuja tecnológica, cuando el Nasdaq se movía por su cuenta. \\
$\tau$ móvil: máx. & 0,832 (2009-09) & Máxima dependencia al final de la crisis financiera. \\
Mediana & 0,751 & Valor típico de la dependencia en ventanas de un año. \\
\end{xltabular}}
